\subsection{Variable selection}

\paragraph{Definitions.} 
\citet{borio2014financial} defines the financial cycle as ``self-reinforcing interactions between perceptions of value and risk, attitudes towards risk, and financing constraints, which translate into booms followed by busts." This definition highlights that the financial cycle is characterised by joint fluctuations in credit, asset prices, and leverage that can become disconnected from underlying fundamentals. As financial imbalances accumulate over time, cyclical systemic risk increases, raising the likelihood of severe downturns and banking crises.

Measuring the financial cycle therefore requires selecting variables that capture both the build-up of financial imbalances and their potential materialisation. In this paper, we focus on four key dimensions: credit, asset prices, indebtedness, and financial conditions. These dimensions are chosen based on their established role in the literature as core components of cyclical systemic risk and their relevance for early-warning analysis.

To preserve interpretability and ensure a transparent link between the estimated factor and its underlying drivers, we deliberately restrict the analysis to a parsimonious set of variables within each dimension. This approach facilitates the economic interpretation of the financial cycle indicator while retaining the key information needed to capture its dynamics.

\paragraph{Credit.}  

The central role of credit in the financial cycle is well established in both theoretical and empirical work. Early contributions by \citet{fisher1933debt} and \citet{minsky1977financial} highlight the role of credit cycles in generating and amplifying macroeconomic instability, while subsequent work shows how credit acts as a propagation mechanism for shocks \citep{bernanke1983non, gertler1988financial}. This mechanism is formalised in the financial accelerator framework of \citet{bernanke1999financial}, which emphasises the endogenous amplification of shocks through credit market conditions. Empirically, credit aggregates have been shown to be among the most reliable predictors of systemic banking crises \citep{schularick2012credit}.

To capture these dynamics, we include two measures of credit: credit to non-financial corporations and credit to households (including non-profit institutions). This distinction allows us to account for differences in borrowing behaviour across sectors, which may have distinct implications for the build-up of cyclical systemic risk. Both series are constructed from the total stock of credit to the private non-financial sector, covering all sources of financing including domestic banks, non-bank financial institutions, non-financial corporations, and cross-border lending. As such, the indicator provides a comprehensive measure of credit developments in the economy.

\paragraph{Asset prices.}  

Asset prices play a central role in the financial cycle, both as indicators of financial imbalances and as channels through which shocks are amplified. In particular, increases in asset prices can relax borrowing constraints and support credit expansion, while subsequent declines can trigger adverse feedback loops by reducing collateral values \citep{kiyotaki1997credit}. Empirically, asset price dynamics—especially in housing markets—are closely associated with the build-up and materialisation of systemic risk \citep{jorda2015leveraged, claessens2012business}. House prices, in particular, have been shown to be strong predictors of banking crises, while equity prices often peak prior to crisis episodes \citep{mian2014explains, reinhart2013banking}.

To capture these dynamics, we include both house price and share price indices. The house price index reflects inflation-adjusted transaction prices of residential property, providing a measure of developments in real estate markets that are closely linked to credit conditions. Share price indices capture the prices of publicly traded equities, aggregated from daily closing values, and provide forward-looking information on market valuations and risk sentiment.

\paragraph{Indebtedness.} 

While credit captures the expansion of financing over the financial cycle, indebtedness reflects the accumulation of debt and the associated vulnerabilities in the economy. High levels of leverage increase the sensitivity of households, firms, and financial intermediaries to adverse shocks, as declines in income or asset prices can impair debt servicing capacity and amplify financial distress \citep{adrian2010liquidity, brunnermeier2009market}. As a result, elevated indebtedness is a key determinant of the severity of downturns and the materialisation of systemic risk.

To capture these dynamics, we include two complementary indicators: the credit-to-GDP ratio and the debt service ratio (DSR). The credit-to-GDP ratio provides a measure of the stock of debt relative to the size of the economy and is widely used as an indicator of systemic risk, with well-documented early-warning properties for banking crises \citep{aikman2015curbing}. The DSR captures the share of income devoted to debt servicing, providing a direct measure of repayment burden and financial stress faced by households and non-financial corporations. Together, these variables capture both the level and the sustainability of indebtedness over the financial cycle.

\paragraph{Financial conditions.}

Financial conditions capture the pricing of risk and the level of risk appetite in financial markets, both of which are central to the financial cycle. We proxy these conditions using the spread between each country’s 10-year sovereign bond yield and that of Germany, which serves as a benchmark reflecting low credit risk within the European context. Changes in sovereign spreads provide information on shifts in perceived risk and financing conditions, as increases in risk premia tend to coincide with periods of financial stress.

This measure is closely related to indicators used in financial stress indices \citep{hubrich2015financial}, which distinguish periods of low risk appetite from episodes of heightened uncertainty. Moreover, sovereign bond spreads are informative about broader financial conditions, as movements in sovereign risk premia are often mirrored in borrowing costs for households, firms, and financial intermediaries. Including this variable therefore enhances the ability of the indicator to capture the materialisation of systemic risk, complementing other variables that primarily reflect its build-up.

\subsection{Variable treatment}

\paragraph{Stationarity.} 

An underlying assumption of dynamic factor models is that the data are covariance stationary. To satisfy this requirement, all variables, except the government bond spread, are transformed into 2-year growth rates, while the bond spread is expressed in 2-year differences. These transformations are guided by unit root tests and case-by-case judgement.

The choice of 2-year growth rates reflects the medium-term nature of the financial cycle and is supported by evidence that longer-horizon transformations exhibit stronger early-warning properties for systemic banking crises \citep{lang2019anticipating}. By focusing on medium-term dynamics, this approach captures the gradual build-up of financial imbalances that precede crisis episodes.

Importantly, this methodology analyses cycles in growth rates rather than relying on gap measures, such as the credit-to-GDP gap, which extract cyclical components from time series levels. In line with \citet{schuler2020financial} and \citet{lang2019anticipating}, this avoids the spurious cycles and endpoint biases often associated with filtering techniques \citep{schuler2018detrending}, providing a more robust representation of financial cycle dynamics.

\paragraph{Low-Frequency Movements.} 

To isolate the cyclical components of the variables, low-frequency movements are removed from the data. Long-term trends often reflect structural changes and can obscure medium-term fluctuations relevant for the financial cycle. As emphasized by \citet{stock2016dynamic}, removing these low-frequency components is essential for distinguishing persistent trends from cyclical dynamics in macroeconomic time series.

This is particularly relevant in the European context, where greater financial integration and the introduction of the euro have significantly influenced long-term financial dynamics and contributed to the synchronization of financial cycles across countries \citep{jorda2018global, aikman2015curbing}. In addition, the post-GFC period has been marked by significant institutional changes that have reshaped the functioning of financial markets. Accounting for these low-frequency movements is therefore important to avoid conflating structural trends with cyclical fluctuations.

In line with this approach, we apply a filtering method to extract medium-term dynamics \citep{drehmann2012characterising, schuler2020financial}. Specifically, we use the Butterworth filter proposed by \citet{pollock2000trend}, which offers greater flexibility than the Hodrick–Prescott (HP) filter in capturing time-varying trends.

We focus on cycles with periodicities between 4 and 25 years. Frequencies longer than 25 years are removed to exclude low-frequency trends, while fluctuations with periodicities shorter than 4 years are filtered out to reduce high-frequency noise. This approach allows us to isolate medium-term financial cycle dynamics in a consistent and robust manner.

\paragraph{Standardization.}  

All variables are standardized to have zero mean and unit variance. This is a standard step in dynamic factor models, as the estimated factor is sensitive to the scale of the variables. Standardization ensures comparability across variables and allows the resulting financial cycle indicator to be interpreted in terms of deviations from its historical average. In this framework, large positive or negative values of the indicator reflect significant departures from typical conditions and can be interpreted as signals of financial imbalances.

\subsection{Descriptive statistics}

\paragraph{Sample.} Our sample consists of quarterly data for 10 European countries: Sweden (SE), Italy (IT), France (FR), Germany (DE), Belgium (BE), the Netherlands (NL), Finland (FI), the United Kingdom (GB), Portugal (PT), and Spain (ES). The dataset spans from 1995 Q3 to 2024 Q3. The country selection balances sample length with data availability across countries. Our focus on European countries reflects their exposure to similar economic shocks and similar institutional frameworks, enabling more meaningful comparisons. In the following sections, we construct the financial cycle indicator for each country and present the results by examining the distribution of the indicator across the sample.

\paragraph{Descriptive statistics.}  
  
Table \ref{sum_st} reports summary statistics for the selected variables, expressed in 2-year growth rates, except for bond spreads, which are presented in 2-year differences. The results reveal substantial heterogeneity across countries and across types of financial indicators.

Credit to non-financial corporations shows moderate average growth across countries, ranging from 7.3\% in Italy to 14.4\% in Spain, with relatively limited dispersion. In contrast, credit to households grows more strongly and is more volatile, particularly in Portugal (17.2\%), Spain (13.9\%), and Sweden (13.7\%).

Asset price dynamics display greater cross-country variation. House prices record positive average growth in most countries, with higher values in Sweden (9.2\%) and the United Kingdom (8.1\%), while Italy stands out with a negative average growth rate (-0.5\%). Equity prices exhibit robust growth across all countries, reaching 25.5\% in Finland and 22.6\% in Sweden, but with substantial volatility.

Indicators of indebtedness show more moderate dynamics. Debt service ratios display near-zero average growth across countries, ranging from -1.7\% in Germany to 2.1\% in Belgium. By contrast, credit-to-GDP ratios increase more noticeably in countries such as Spain (5.0\%), Belgium (4.3\%), Portugal (3.8\%), and Sweden (4.2\%), indicating a gradual accumulation of leverage over the sample period.

Finally, sovereign bond spreads generally decline over the sample period, particularly in Portugal (-21\%), Sweden (-18\%), and Spain (-17\%). As Germany serves as the benchmark for the spread measure, it is not directly comparable to other countries in this dimension. Variability in spreads is highest in Portugal, reflecting the pronounced fluctuations observed during the sovereign debt crisis.

\begin{table}[htbp]
\centering
\caption{Summary statistics}
\label{sum_st}

\resizebox{\textwidth}{!}{%
\begin{tabular}{l *{14}{r}}
\toprule
& \multicolumn{2}{c}{\shortstack{Credit\\NFC}} 
& \multicolumn{2}{c}{\shortstack{Credit\\Households}} 
& \multicolumn{2}{c}{\shortstack{House\\prices}} 
& \multicolumn{2}{c}{\shortstack{Share\\prices}} 
& \multicolumn{2}{c}{DSR} 
& \multicolumn{2}{c}{\shortstack{Credit-to-\\GDP}} 
& \multicolumn{2}{c}{\shortstack{Bond\\spread}} \\
\cmidrule(lr){2-3}\cmidrule(lr){4-5}\cmidrule(lr){6-7}\cmidrule(lr){8-9}\cmidrule(lr){10-11}\cmidrule(lr){12-13}\cmidrule(lr){14-15}
\textbf{Country} 
& \textbf{mean} & \textbf{sd}
& \textbf{mean} & \textbf{sd}
& \textbf{mean} & \textbf{sd}
& \textbf{mean} & \textbf{sd}
& \textbf{mean} & \textbf{sd}
& \textbf{mean} & \textbf{sd}
& \textbf{mean} & \textbf{sd} \\
\midrule
BE & 13.1\% & 11.1\% & 11.1\% & 5.1\%  & 5.0\%  & 5.4\%  & 13.4\% & 33.4\% & 2.1\%  & 7.8\%  & 4.3\% & 7.9\%  & -0.02 & 0.48 \\
DE & 7.4\%  & 6.8\%  & 5.4\%  & 5.3\%  & 1.1\%  & 7.4\%  & 13.0\% & 30.9\% & -1.7\% & 4.0\%  & 0.7\% & 5.1\%  & -0.32 & 1.07 \\
ES & 14.4\% & 21.2\% & 13.9\% & 19.3\% & 4.9\%  & 14.3\% & 13.7\% & 38.5\% & 0.3\%  & 13.6\% & 5.0\% & 13.6\% & -0.17 & 1.33 \\
FI & 9.6\%  & 8.2\%  & 12.3\% & 9.6\%  & 3.7\%  & 8.2\%  & 25.5\% & 63.8\% & 0.1\%  & 10.4\% & 2.2\% & 7.2\%  & -0.11 & 0.50 \\
FR & 9.3\%  & 4.7\%  & 11.0\% & 5.4\%  & 5.6\%  & 10.2\% & 14.4\% & 30.9\% & 0.9\%  & 5.9\%  & 3.5\% & 4.7\%  & 0.01  & 0.31 \\
GB & 10.7\% & 10.8\% & 11.4\% & 8.1\%  & 8.1\%  & 11.8\% & 8.5\%  & 21.4\% & 0.2\%  & 7.4\%  & 2.3\% & 7.6\%  & 0.03  & 0.52 \\
IT & 7.3\%  & 8.7\%  & 12.1\% & 10.8\% & -0.5\% & 8.6\%  & 12.4\% & 36.9\% & 0.1\%  & 8.8\%  & 3.0\% & 7.9\%  & -0.16 & 1.46 \\
NL & 10.1\% & 5.0\%  & 12.8\% & 10.2\% & 7.2\%  & 11.5\% & 14.9\% & 33.8\% & 0.1\%  & 5.6\%  & 2.4\% & 5.8\%  & 0.02  & 0.17 \\
PT & 10.7\% & 11.7\% & 17.2\% & 22.3\% & 2.4\%  & 9.4\%  & 17.8\% & 39.5\% & -0.6\% & 9.4\%  & 3.8\% & 10.8\% & -0.21 & 2.68 \\
SE & 13.5\% & 10.3\% & 13.7\% & 5.6\%  & 9.2\%  & 9.4\%  & 22.6\% & 34.7\% & 1.5\%  & 8.1\%  & 4.2\% & 7.9\%  & -0.18 & 0.69 \\
\bottomrule
\end{tabular}%
}

\vspace{0.3cm}
\footnotesize
\parbox{0.95\textwidth}{\textbf{Notes}: NFC refers to non-financial corporations. Bond spreads are defined as the difference between the 10-year sovereign bond yield of each country and that of Germany (DE), which serves as the benchmark. For Germany, the reported series corresponds to the 10-year sovereign bond yield. Summary statistics are computed using 2-year growth rates for all variables, except for bond spreads, which are expressed as 2-year differences. The abbreviation sd denotes the standard deviation.}
\end{table}
